\section{Server API}
\label{sec:ServerApi}

This chapter documents the WebSocket and REST API endpoints exposed by the
Archimesh server.

\subsection{WebSocket endpoint}

\begin{description}
\item[Path] \code{/models/\{modelId\}/stream}
\item[Protocol] WebSocket (RFC 6455)
\item[Sub-protocol] None (plain JSON text frames)
\end{description}

The WebSocket endpoint handles real-time collaboration. Each connection is
scoped to one model.

\subsubsection{WebSocket lifecycle}

\begin{enumerate}
\item Client opens a WebSocket connection to the endpoint.
\item Server performs authorization checks (see
  Section~\ref{sec:Authorization}) and either accepts the connection or
  closes it with policy violation status.
\item Client sends a \code{Join} message to request model state.
\item Server responds with \code{CheckoutSnapshot} or
  \code{CheckoutDelta}.
\item Client sends \code{SubmitOps}, \code{AcquireLock},
  \code{ReleaseLock} and \code{Presence} messages during the session.
\item Server replies to the submitting session with \code{OpsAccepted}
  or a correlated \code{Error}. Model subscribers receive
  \code{OpsBroadcast}, \code{LockEvent} and \code{PresenceBroadcast}.
\item Client or server closes the WebSocket. On close the server unregisters
  the session and releases held locks.
\end{enumerate}

\subsubsection{Inbound messages (client to server)}

All inbound messages are wrapped in a \code{ClientEnvelope}:

\begin{lstlisting}[caption={ClientEnvelope shape}]
{
  "type": "<MessageType>",
  "payload": { ... }
}
\end{lstlisting}

\begin{longtable}{p{32mm}p{52mm}p{70mm}}
\toprule
Type & Purpose & Key payload fields \\
\midrule
\code{Join} & Request model checkout &
  \code{lastSeenRevision} (optional),
  \code{ref} (optional; \code{HEAD} or tag name),
  \code{actor} (\code{userId}, \code{sessionId}) \\
\code{SubmitOps} & Submit one batch of operations &
  \code{baseRevision},
  \code{opBatchId},
  \code{actor},
  \code{ops} (array) \\
\code{AcquireLock} & Acquire a notation lease &
  \code{actor},
  \code{targets} (array of entity ids),
  \code{ttlMs} \\
\code{ReleaseLock} & Release a notation lease &
  \code{actor},
  \code{targets} (array of entity ids) \\
\code{Presence} & Broadcast cursor/selection &
  \code{actor},
  \code{viewId},
  \code{selection} (array),
  \code{cursor} (object) \\
\bottomrule
\end{longtable}

\subsubsection{Join and checkout}

When the server receives a \code{Join}:

\begin{enumerate}
\item The model must exist (registered via admin API).
\item Authorization is checked (join action).
\item The requested ref is resolved (\code{HEAD} or tag name).
\item Tagged refs always receive their stored read-only snapshot. For HEAD,
  a positive trusted client revision at or below head can receive a delta
  only when every range and operation count in the requested window is present.
  A gap after pruning falls back to a current snapshot.
\item Otherwise, the server responds with \code{CheckoutSnapshot}
  containing the full materialized model state at the requested ref.
\end{enumerate}

\subsubsection{SubmitOps processing}

The \code{SubmitOps} pipeline (implemented in \code{ArchimeshService}):

\begin{enumerate}
\item Model registration and writable ref checks.
\item Validation of the op batch structure.
\item Normalization and expansion (bootstrap view ref normalization,
  cascade delete expansion).
\item Durable idempotency check: an existing batch receipt returns its original
  \code{OpsAccepted} range, including after payload pruning.
\item Revision-ahead check: if \code{baseRevision > headRevision},
  send \code{Error(REVISION\_AHEAD)} to the submitting session.
\item Precondition validation: all entity references must exist either in
  persisted state or in earlier ops within the same batch.
\item Lock conflict check: for operations requiring notation locks
  (\code{UpdateViewObjectOpaque}, \code{UpdateConnectionOpaque}),
  verify no other session holds a conflicting lock.
\item Revision range assignment: allocate \code{from..to} consuming one
  revision per op.
\item Causal metadata normalization: stamp each op with
  \code{clientId}, \code{lamport}, \code{opId}.
\item Atomic persistence: append log, apply projection, update head and store
  pending publication payload in one Neo4j transaction. Failed persistence
  receives no acceptance and invalidates the in-memory revision cache.
\item Cache the committed acceptance range.
\item Optional post-commit consistency checks.
\item Request the durable delivery worker to publish committed batches.
\item Reply \code{OpsAccepted} to the submitting session.
\end{enumerate}

If validation rejects the batch, a correlated \code{Error} is sent to its submitter
with a specific code (\code{REVISION\_AHEAD}, \code{PRECONDITION\_FAILED},
\code{LOCK\_CONFLICT}).

\subsubsection{Outbound messages (server to client)}

\begin{longtable}{p{36mm}p{52mm}p{66mm}}
\toprule
Type & Purpose & Key payload fields \\
\midrule
\code{CheckoutSnapshot} & Full model state &
  \code{modelId}, \code{headRevision},
  \code{elements}, \code{relationships},
  \code{views}, \code{folders}, etc. \\
\code{CheckoutDelta} & Incremental changes &
  \code{modelId},
  \code{fromRevision}, \code{toRevision},
  \code{ops} (array) \\
\code{OpsAccepted} & Batch commit confirmation &
  \code{opBatchId}, \code{baseRevision},
  \code{assignedRevisionRange} \\
\code{OpsBroadcast} & Remote batch broadcast &
  \code{opBatch} (full committed batch) \\
\code{LockEvent} & Lock state change &
  \code{action} (acquired/released/expired),
  \code{actor}, \code{targets} \\
\code{PresenceBroadcast} & Remote cursor/selection &
  \code{actor}, \code{viewId},
  \code{selection}, \code{cursor} \\
\code{Error} & Rejection or failure &
  \code{code}, \code{message}, optional \code{opBatchId} \\
\bottomrule
\end{longtable}

\subsection{REST API}

The server exposes REST endpoints under the base path \code{/}.

\subsubsection{Model state endpoints}

\begin{longtable}{p{54mm}p{44mm}p{56mm}}
\toprule
Endpoint & Method & Purpose \\
\midrule
\code{/models/\{modelId\}/snapshot} & GET & Get the current materialized
  snapshot of a model. Optional query parameter \code{?ref=<tagName>}. \\
\code{/models/\{modelId\}/rebuild} & POST & Rebuild materialized state
  from checkpoint plus complete tail, or full history. Transactional replacement. \\
\bottomrule
\end{longtable}

Readable-model discovery uses \code{GET /models} and
\code{GET /models/\{modelId\}/tags}. These filter by read access; ordinary
plugin users do not need access to the admin catalog.

Example:

\begin{lstlisting}[language=bash,caption={Get snapshot}]
curl -s http://localhost:8081/models/demo/snapshot | jq .
\end{lstlisting}

\subsubsection{Admin endpoints}

\begin{longtable}{p{54mm}p{44mm}p{56mm}}
\toprule
Endpoint & Method & Purpose \\
\midrule
\code{/admin/models} & GET & List registered models (catalog). \\
\code{/admin/models/\{modelId\}} & POST & Register a new model. \\
& PUT & Rename with query parameter \code{modelName}. \\
& DELETE & Delete a model. Query param
  \code{?force=true} to override active-session guard. \\
\code{/admin/models/\{modelId\}/status} & GET & Aggregated status:
  snapshot counts, revision, consistency. \\
\code{/admin/models/\{modelId\}/rebuild-and-status} & POST & Rebuild and
  return status in one call. \\
\code{/admin/models/\{modelId\}/window?limit=25} & GET & Model window:
  status + recent activity + recent op batches. \\
\code{/admin/overview?limit=25} & GET & Overview for all known
  active/recent models. \\
\code{/admin/models/\{modelId\}/integrity} & GET & Integrity report:
  missing references, orphan nodes. \\
\code{/admin/models/\{modelId\}/acl} & GET, PUT & Read or update
  per-model ACL (admin, writer, reader user lists). \\
\code{/admin/models/\{modelId\}/tags} & POST & Create an immutable tag
  at the current HEAD. \\
\code{/admin/models/\{modelId\}/export} & GET & Export model package
  (snapshot, recovery data, tags and ACL). \\
\code{/admin/models/import?overwrite=true|false} & POST & Import a model
  package. \\
\code{/admin/models/\{modelId\}/checkpoints} & POST & Create immutable checkpoint. \\
\code{/admin/models/\{modelId\}/checkpoints/latest} & GET & Latest checkpoint summary. \\
\code{/admin/models/\{modelId\}/compact} & POST & Run compaction on
  the model (prune covered published payloads; supports dryRun). \\
\code{/admin/auth/diagnostics} & GET & Authorization diagnostics:
  resolved identity, subject, normalized roles. \\
\code{/admin/sessions} & GET & Active WebSocket sessions across all
  models. \\
\code{/admin/audit/config} & GET & Current audit configuration. \\
\bottomrule
\end{longtable}

Checkpoint creation and inspection are model-admin actions. The latest
summary returns model ID, revision, creation timestamp, SHA-256 hash and raw
node/edge counts; no checkpoint gives 404. Compaction accepts
\code{retainRevisions} and \code{dryRun=true|false} query parameters.
Its report includes the verified checkpoint boundary, eligibility before the
run, actual deletions and durable pruning progress. See
Section~\ref{sec:Management} for the full maintenance workflow.

Examples:

\begin{lstlisting}[language=bash,caption={Admin API examples}]
# List all models
curl -s http://localhost:8081/admin/models | jq .

# Create a model
curl -s -X POST http://localhost:8081/admin/models/demo | jq .

# Get model status
curl -s http://localhost:8081/admin/models/demo/status | jq .

# Get model window
curl -s http://localhost:8081/admin/models/demo/window?limit=25 | jq .

# Get overview
curl -s http://localhost:8081/admin/overview?limit=25 | jq .

# Delete model
curl -s -X DELETE http://localhost:8081/admin/models/demo | jq .

# Force delete (override active-session guard)
curl -s -X DELETE "http://localhost:8081/admin/models/demo?force=true" | jq .
\end{lstlisting}

\subsubsection{Admin UI endpoint}

\begin{description}
\item[Path] \code{/admin-ui/}
\item[Serves] Static Svelte admin application (built from
  \filepath{admin-web/}).
\end{description}

When running in dev mode, the admin UI is served from
\filepath{server/src/main/resources/META-INF/resources/admin-ui/} (populated
by the build script). In production, the Quarkus application serves the
built static assets.
